%=======================================================================
%  CSE558 Data Science - End-term Project Report (built on the official template, same layout as the midterm)
%  Monsoon 2026, IIIT-Delhi.  Group DSC 10.
%  Every number comes from 3_endsem/results/*.csv or *.json (see README, "number -> file" table).
%  Compile:  tectonic main.tex
%=======================================================================
\documentclass[10pt,a4paper,twocolumn]{article}

\newcommand{\projecttitle}{Smaller data, same forecast: sensitivity-sampling coresets for 42-day hub demand,
  and an audited, frozen forecaster}

%-------------------------------------------- fonts (book serif body, every measured quantity in mono)
\usepackage{iftex}
\ifPDFTeX
  \usepackage[T1]{fontenc}\usepackage[utf8]{inputenc}\usepackage{XCharter}\usepackage[scaled=0.86]{beramono}
\else
  \usepackage{fontspec}
  \setmainfont{Charter}[BoldFont={Charter Bold},ItalicFont={Charter Italic},BoldItalicFont={Charter Bold Italic}]
  \setsansfont{Avenir Next}[Scale=0.92]
  \setmonofont{Menlo}[Scale=0.80]
\fi

%----------------------------------------------------------- packages
\usepackage[margin=16mm,bottom=19mm,columnsep=6.5mm]{geometry}
\usepackage{amsmath,amssymb}
\usepackage{xcolor}
\usepackage{graphicx}
\usepackage{tabularx,array,booktabs,multirow}
\usepackage{enumitem}
\usepackage{titlesec}
\usepackage{needspace}
\usepackage[font=small,labelfont={bf,sf}]{caption}
\usepackage[most]{tcolorbox}
\usepackage{fancyhdr}
\usepackage{lastpage}
\usepackage{tikz}
\usepackage{hyperref}

%------------------------------------------------------------- tokens (DIP "photometric" series)
\definecolor{ink}{HTML}{0E1116}
\definecolor{greytext}{HTML}{4A515E}
\definecolor{greysoft}{HTML}{878E9A}
\definecolor{rule}{HTML}{D8DCE2}
\definecolor{ultra}{HTML}{1B36C4}     % only: exactly verified claims, our method
\definecolor{amber}{HTML}{C77800}     % only: limitations, losses
\definecolor{wa}{HTML}{1F2430}\definecolor{wb}{HTML}{4A515E}\definecolor{wc}{HTML}{878E9A}
\definecolor{wd}{HTML}{C4C9D0}\definecolor{we}{HTML}{E8EBEF}
\hypersetup{colorlinks=true, urlcolor=ultra, linkcolor=ultra, citecolor=ultra, breaklinks=true,
            pdftitle={\projecttitle}, pdfsubject={CSE558 Data Science, Monsoon 2026: End-term Project Report}}
\color{ink}

\newcommand{\stepwedge}{\par\noindent{\setlength{\fboxsep}{0pt}%
  \colorbox{wa}{\makebox[0.2\linewidth]{\rule{0pt}{1.4mm}}}\colorbox{wb}{\makebox[0.2\linewidth]{}}%
  \colorbox{wc}{\makebox[0.2\linewidth]{}}\colorbox{wd}{\makebox[0.2\linewidth]{}}\colorbox{we}{\makebox[0.2\linewidth]{}}}\par}
\newcounter{glyph}
\newcommand{\glyph}{\stepcounter{glyph}\begin{tikzpicture}[x=1.35mm,y=1.35mm,baseline=0.2mm]
  \foreach \i in {0,...,4}{\foreach \j in {0,...,4}{\fill[greysoft!55] (\i,-\j) circle (0.18);}}
  \pgfmathtruncatemacro{\gi}{mod(\value{glyph}-1,5)}\pgfmathtruncatemacro{\gj}{div(\value{glyph}-1,5)}
  \fill[ultra] (\gi,-\gj) circle (0.36);\end{tikzpicture}}
\titleformat{\section}[display]{\normalfont\large\bfseries}{\glyph\hspace{1.6mm}{\sffamily\footnotesize\color{greytext}SECTION~\thesection}}{1pt}{}
\titlespacing*{\section}{0pt}{8pt}{2pt}
\titleformat{\subsection}{\normalfont\normalsize\bfseries}{\thesubsection}{0.5em}{}
\titlespacing*{\subsection}{0pt}{6pt}{1pt}
\newcommand{\sectionbreak}{\needspace{8\baselineskip}}

\setlength{\parindent}{0pt}
\setlength{\parskip}{3.2pt}
\setlist{leftmargin=1.1em, itemsep=1pt, topsep=1pt, parsep=0pt}
\renewcommand{\arraystretch}{1.12}
\setlength{\tabcolsep}{3.2pt}
\newcolumntype{L}[1]{>{\raggedright\arraybackslash}p{#1}}
\newcolumntype{Y}{>{\raggedright\arraybackslash}X}
\newcolumntype{M}{>{\ttfamily}r}
\newcolumntype{O}{>{\ttfamily\bfseries\color{ultra}}r}
\newcommand{\key}[1]{\textbf{\textcolor{ultra}{#1}}}
\newcommand{\warn}[1]{\textbf{\textcolor{amber}{#1}}}
\newcommand{\n}[1]{\texttt{#1}}
\newcommand{\code}[1]{\texttt{#1}}
\newcommand{\thh}[1]{{\sffamily\scriptsize\bfseries\color{greytext}\MakeUppercase{#1}}}
\newcommand{\FIG}{../figures/}

\newtcolorbox{stamp}[1]{enhanced,colback=white,colframe=ultra,boxrule=0.6pt,arc=0pt,left=3pt,right=3pt,top=2pt,bottom=2pt,
  fonttitle=\sffamily\scriptsize\bfseries,coltitle=ultra,colbacktitle=white,title=\MakeUppercase{#1},
  attach title to upper,after title={\par\smallskip},fontupper=\footnotesize}
\newtcolorbox{limit}[1]{enhanced,colback=amber!6,colframe=amber,boxrule=0pt,leftrule=1.6pt,arc=0pt,left=4pt,right=3pt,top=2pt,bottom=2pt,
  fonttitle=\sffamily\scriptsize\bfseries,coltitle=amber,title=\MakeUppercase{#1},attach title to upper,after title={\par},fontupper=\footnotesize}
\newtcolorbox{eqbox}{enhanced,colback=ultra!3,colframe=ultra,boxrule=0pt,leftrule=1.6pt,arc=0pt,left=4pt,right=3pt,top=1pt,bottom=1pt}

\newcommand{\readout}[2]{\begin{minipage}[t]{0.19\linewidth}{\ttfamily\large\color{ultra}#1}\\[-1pt]{\sffamily\scriptsize\color{greytext}#2}\end{minipage}}

\pagestyle{fancy}\fancyhf{}\renewcommand{\headrulewidth}{0pt}\renewcommand{\footrulewidth}{0.3pt}
\fancyfoot[L]{\sffamily\scriptsize\color{greytext}CSE558 \textbullet\ DSC 10 \textbullet\ End-term report \textbullet\ Shadow Hub}
\fancyfoot[R]{\ttfamily\scriptsize\color{greytext}\thepage\ / \pageref{LastPage}}
\fancypagestyle{plain}{\fancyhf{}\fancyfoot[L]{\sffamily\scriptsize\color{greytext}CSE558 \textbullet\ DSC 10 \textbullet\ End-term report}\fancyfoot[R]{\ttfamily\scriptsize\color{greytext}\thepage\ / \pageref{LastPage}}}

%=======================================================================
\begin{document}

\twocolumn[%
{\sffamily\footnotesize\bfseries\color{ultra}CSE558 DATA SCIENCE \ \textbullet\ \ END-TERM PROJECT REPORT \ \textbullet\ \ MONSOON 2026\par}
\vspace{3pt}
{\LARGE\bfseries\raggedright\projecttitle\par}
\vspace{3pt}
{\normalsize\itshape\color{greytext} A 50{,}000-row sensitivity coreset (4.4\% of the training rows) brings LightGBM within 0.0048 RMSLE of the
full 1.13-million-row fit and beats uniform sampling at every mid size. Separately, the forecaster we submit (tuned LightGBM blended
with TabPFN-TS) improves on the same LightGBM recipe by 0.0045 on development folds and was frozen only after an independent
re-computation of its adoption gate.\par}
\vspace{4pt}
\readout{0.1022}{full LightGBM, test RMSLE}\hfill
\readout{0.1070}{coreset m = 50k (uniform 0.1102)}\hfill
\readout{2.75$\times$}{speed-up incl. build}\hfill
\readout{2.47}{total sensitivity S (in [2,6])}\hfill
\readout{+0.0045}{submitted forecaster, gain folds 2--4}
\vspace{4pt}
\stepwedge
\vspace{3pt}
{\small
\begin{tabularx}{\linewidth}{@{}L{2.0cm} Y@{}}
\toprule
Group & DSC 10 \quad\textbullet\quad Dataset: Shadow Hub database, 1{,}115 hubs, 1 Jan 2013 -- 19 Jun 2015 (raw files in
  \url{https://github.com/kushiiitd05/ShadowDsc10-ShadowHub/tree/main/1_midterm/data/raw}) \\
Code \& README & \url{https://github.com/kushiiitd05/ShadowDsc10-ShadowHub/tree/main/3_endsem}
  (notebooks \code{notebook/01\_endterm\_results.ipynb} and \code{02\_run\_frozen\_forecaster.ipynb}, scripts in \code{code/},
  frozen models in \code{final/}) \\
\bottomrule
\end{tabularx}}
\vspace{2pt}
{\small
\begin{tabularx}{\linewidth}{@{}L{0.4cm} L{3.1cm} L{1.6cm} L{4.3cm} Y@{}}
\toprule
\thh{\#} & \thh{Name} & \thh{Roll no.} & \thh{Institute email} & \thh{End-term contribution} \\
\midrule
1 & Manya Chaturvedi & 2023312 & manya23312@iiitd.ac.in & \warn{to be confirmed by the member} \\
2 & Kush Tokas & 2023296 & kush23296@iiitd.ac.in & model track (experiments, searches, TabPFN-TS blend, gate, freeze audit,
  final retrain); end-term coreset runs; report, notebooks, repository \\
3 & Akash Adur & 2024053 & akash24053@iiitd.ac.in & \warn{to be confirmed by the member} \\
4 & Anusha Rana & 2023114 & anusha23114@iiitd.ac.in & \warn{to be confirmed by the member} \\
5 & Sanjeev & 2023483 & sanjeev23483@iiitd.ac.in & \warn{to be confirmed by the member} \\
\bottomrule
\end{tabularx}}
\vspace{8pt}
]

%-----------------------------------------------------------------------
\section{Problem, data and the fixed split}

\begin{figure*}[t]
\centering\includegraphics[width=0.94\linewidth]{\FIG fig_pipeline.pdf}
\caption{\textsf{PIPELINE.} One feature pipeline, two tracks. Track A answers the course rubric with $x^\ast$ = LightGBM
\code{cs02} on a fixed split. Track B is the forecaster we submit; it is improved on rolling development folds and frozen.
Track B is not coreset-ed (box in Section~1).}
\end{figure*}

\textbf{Task.} Forecast the daily order volume of each of 1{,}115 retail hubs for the next 42 days. One row is one hub on one
day. The response is skewed (median about 6{,}000 on open days, 0 on closed days), so every model is fitted on
$\log(1+y)$ and scored by RMSLE $=\sqrt{\tfrac1N\sum(\log(1+\hat y)-\log(1+y))^2}$; RMSE on the log scale and $R^2$ on the
log scale are reported beside it, as the end-term brief asks.

\begin{stamp}{Reading guide: two tracks, three model names}
\textbf{Track A, the rubric (Sections 2--5).} $x^\ast$ = LightGBM recipe \code{cs02} (63 leaves, feature fraction 0.2, row
bagging), trained on the fixed split; ridge is the second model; every summary is fitted with the same settings.
\textbf{Track B, the forecaster we submit (Sections 6--7).} \code{hp17} (the same LightGBM, tuned: 255 leaves, feature
fraction 0.15, strong L2), 3 seeds, blended 0.8/0.2 in log space with TabPFN-TS (365- and 730-day contexts, 0.10 each). It
writes the forecast file for the unlabelled window from 20 Jun 2015, so it has development-fold scores only.
\textbf{Why Track B is not coreset-ed.} TabPFN-TS is a pretrained model that reads each hub's history at prediction time:
it has no training loss $F$, no fitted $x^\ast$ and no rows to weight, so sensitivities are not defined for it; and the rubric's
same-model rule would need the whole blend refitted for 5 repeats $\times$ 7 sizes $\times$ 3 methods.
\end{stamp}

\textbf{Rows and columns.} Each training row is a (forecast origin, hub, target day) triple: for an origin $o$ the features
are built from history up to $o$ only, and the label is the volume on day $o+h$, $h=1..42$. We build 37 origins two weeks
apart (label windows end by 19 Jun 2015) with \n{342} engineered columns: calendar, promotion and holiday flags known in
advance, hub metadata, rolling medians over 7--365 days, same-weekday means, a cross-fitted hub target encoding, and data-gap
flags. The end-term training split has \n{1,133,640} open rows with positive volume; $1{,}133{,}640\times342=\n{387.7}$
million cells, far above the 3-million floor of the course.

\textbf{Split (fixed once, before any end-term run).} Train = every origin whose 42-day label window ends on or before
8 May 2015; test = the last 42 training days, 9 May -- 19 Jun 2015 (\n{46,830} rows = 1{,}115 hubs $\times$ 42 days).
Hyper-parameters of the full model $x^\ast$ are set on a training-only holdout and then frozen. Closed days are \emph{kept}
in the test set with prediction 0 (their error is 0); this inflates $R^2$, so RMSLE without the open-but-zero rows is also
reported.

\textbf{Data preparation carried over from the midterm, and what changed for forecasting.} The midterm cleaned the raw file
(dates parsed, 180 hubs with a six-month gap in 2014 flagged, open days with zero volume flagged, weekday checked against
the date). For forecasting we added four rules, each because a check failed without it:
\begin{itemize}
\item \emph{Features stop at the origin.} Every rolling statistic for origin $o$ is computed from days $\le o$; a test asserts
  that no history date exceeds the origin. Without it a 7-day median would read the first forecast days.
\item \emph{Training-origin purge.} An origin is used for training only if its whole 42-day label window ends on or before the
  forecast origin; otherwise a training label would sit inside the evaluated window.
\item \emph{Cross-fitted hub encoding.} The hub's mean log volume is a strong feature, but computed on the same rows it is fitted
  on it leaks the label; scikit-learn's \code{TargetEncoder.fit\_transform} uses internal cross-fitting, and the earlier
  \code{fit().transform()} pattern (which does not) was replaced.
\item \emph{Fixed rounds.} Early stopping on the evaluated fold leaks; the final setting uses 1000 rounds at learning rate
  0.025, fixed before scoring.
\end{itemize}

\begin{limit}{Overlap with model selection, and the check we ran}
The end-term test window is the same 42 days as development fold~4, whose scores were part of earlier recipe decisions
(\code{cs02} over \code{base}; the forecaster's gate and weight). No fold-4 label enters any end-term fit, sampling
probability, column choice, number of rounds or ridge $\lambda$. To test whether the overlap mattered we re-took every decision
with fold~4 removed (\code{code/leakage\_check.py}, table below): \textbf{none changes}. The absolute test RMSLE is still a
development-fold number, not an untouched holdout; the coreset-versus-uniform comparison is unaffected, because both use the
identical model and settings.
\end{limit}

{\footnotesize
\begin{tabularx}{\linewidth}{@{}Y L{2.15cm} L{2.6cm}@{}}
\toprule
\thh{Decision} & \thh{with fold 4} & \thh{fold 4 removed} \\
\midrule
$x^\ast$ recipe: \code{cs02} over \code{base} & adopt, +0.0093 & adopt, +0.0119 (fold 1 +0.0210) \\
\code{hp17} alone (gate) & fails, +0.0019 & fails (seed 08 loses fold 2) \\
\code{hp17} + TabPFN-TS, $w$ 0.2 (gate) & passes, +0.0045 & passes, +0.0050, every seed \\
TabPFN-TS weight & 0.2 kept (0.3: +0.0007) & 0.2 kept (0.3: +0.0006) \\
365/730 context mix & kept (730: +0.0008) & kept (730: +0.0010) \\
rounds, ridge $\lambda$, learning rate, sensitivities, columns & fold 4 not used & -- \\
\bottomrule
\end{tabularx}}
\captionof{table}{Leakage check: gains are RMSLE over the baseline; a rival setting must win by more than the noise \n{0.0012} to replace the chosen one.}

%-----------------------------------------------------------------------
\section{The full model \texorpdfstring{$x^\ast$}{x*} and a second model}

\textbf{Why LightGBM.} Volume depends on interactions (weekday $\times$ promotion $\times$ hub format, holiday
proximity, data gaps) that a tree ensemble captures without hand-made products, and LightGBM handles 1.1M rows $\times$ 342
columns in about a minute. Recipe (\code{cs02}): 63 leaves, feature fraction 0.2, row bagging 0.85 every iteration,
min child 80, learning rate 0.05, \n{1184} rounds fixed from the training-only holdout. \textbf{Why ridge as the second
model.} It is the model for which coreset theory is exact (squared loss, closed form), so it shows what the theory
predicts; $\lambda=0.1$ was chosen on the same holdout. Both models are trained with identical settings on the full data
and on every summary.

{\small
\begin{tabularx}{\linewidth}{@{}Y M M M M@{}}
\toprule
\thh{Full model, test split} & \thh{RMSLE} & \thh{excl. open-0} & \thh{$R^2$ (log)} & \thh{fit (s)} \\
\midrule
LightGBM $x^\ast$ & \textcolor{ultra}{\textbf{0.1022}} & 0.0932 & 0.9992 & 59.7 \\
Ridge & 0.1592 & 0.1539 & 0.9980 & 4.3 \\
\bottomrule
\end{tabularx}}

The tree model is 36\% lower in RMSLE than ridge. The very high $R^2$ comes mostly from the closed rows (0 vs 0) and from the
large spread between hubs; the RMSLE is the number to read. Times are from a shared server (Section~4).

\textbf{Why these features.} Every column is known at the forecast origin: history up to the origin, and the calendar,
promotion, opening and holiday flags that the forecast file supplies in advance. The figure below reads the split gain of the
frozen boosters (\code{hp17}, the same 342 columns) by group. The hub $\times$ weekday $\times$ promotion profile (12 columns)
carries \n{46\%} of the gain, the rolling order windows \n{22\%}, the promotion calendar \n{8\%} and the hub encoding \n{6\%};
the 115 session and conversion columns carry \n{2.5\%} and hub metadata \n{0.3\%}. Our reading: a hub's daily volume is mostly
``this hub, on this weekday, with or without a promotion'', scaled by its recent level, and the known-ahead flags adjust single
days. This is also why 40 importance-chosen columns already reach \n{0.1102} (Section~3). Gain shares credit among correlated
columns, so we read groups, not single columns, and it is not a causal statement.

\begin{figure}[h]
\centering\includegraphics[width=\linewidth]{\FIG fig_feature_groups.pdf}
\caption{\textsf{FEATURES.} Share of split gain by feature group (mean of the three frozen boosters, min--max whisker);
\code{results/feature\_group\_importance.csv}.}
\end{figure}

%-----------------------------------------------------------------------
\section{Sampling: rows (Plot 1) and columns}

\textbf{Methods.} (i) \emph{Uniform}: $m$ rows with replacement, weight $n/m$ each (the baseline on every plot).
(ii) \emph{Hub-stratified}: equal rows per hub (added because the smallest size, 2{,}000, is close to the hub count).
(iii) \emph{Sensitivity coreset} (Section~5): $p_i=s_i/S$, weight $1/(m\,p_i)$. Sizes $m\in\{2,5,10,25,50,100,200\}\times10^3$;
each size is repeated 5 times with new seeds for sampling, rough models and training; we plot the median and the min--max band.

\begin{figure}[t]
\centering\includegraphics[width=\linewidth]{\FIG fig_plot1_rows.pdf}
\caption{\textsf{PLOT 1.} Test RMSLE against summary size. Look at LightGBM between 5k and 100k: the coreset curve (blue
diamonds) sits below uniform at every size, by \n{0.0032} at 50k. For ridge all methods lie within \n{0.003} of the full
model from 5k rows on, so sampling barely matters.}
\end{figure}

{\small
\begin{tabularx}{\linewidth}{@{}r M M O M@{}}
\toprule
\thh{m} & \thh{uniform} & \thh{hub-strat.} & \thh{coreset} & \thh{coreset$-$unif.} \\
\midrule
2k   & 0.1297 & 0.1322 & 0.1318 & \textcolor{amber}{+0.0021} \\
5k   & 0.1218 & 0.1221 & 0.1195 & $-$0.0023 \\
10k  & 0.1187 & 0.1203 & 0.1170 & $-$0.0017 \\
25k  & 0.1133 & 0.1131 & 0.1108 & $-$0.0025 \\
50k  & 0.1102 & 0.1117 & 0.1070 & $-$0.0032 \\
100k & 0.1099 & 0.1087 & 0.1058 & $-$0.0041 \\
200k & 0.1055 & 0.1054 & 0.1045 & $-$0.0010 \\
\midrule
full & \multicolumn{4}{r}{\ttfamily 0.1022 (1,133,640 rows)} \\
\bottomrule
\end{tabularx}}
\captionof{table}{LightGBM test RMSLE, median of 5 repeats.}

\textbf{Reading.} The coreset wins from 5k to 200k rows; at 2k it ties or loses (the band overlaps), because a
tree with min child 80 cannot use the extra weight variance at that size. Hub-stratified sampling is no better than
uniform: equal hub coverage is already reached at 5k rows. At 50k (\n{4.4\%} of the rows) the coreset is \n{0.0048} above
the full model; uniform needs between 100k and 200k rows for the same error.

\textbf{Columns (one factor at a time; rows fixed at 200k uniform).} Trees cannot use a random projection of the features
(splits need the original columns), so we compare $k$ columns chosen by gain importance from a rough model fitted on a
uniform training sample against $k$ random columns.

\begin{figure}[h]
\centering\includegraphics[width=0.92\linewidth]{\FIG fig_plot1_columns.pdf}
\caption{\textsf{COLUMNS.} Importance-chosen columns (blue) against random columns (grey). With 40 of 342 columns the
importance subset scores \n{0.1102}, random columns \n{0.1770}.}
\end{figure}

Medians, importance vs random: $k=10$ \n{0.1324} vs \n{0.2316}; $k=20$ \n{0.1204} vs \n{0.1658}; $k=80$ \n{0.1087} vs
\n{0.1248}; $k=160$ \n{0.1058} vs \n{0.1129}. Fewer than half of the columns carry almost all of the signal; 160 random columns
lose \n{0.0071} against 160 chosen ones. (A Gaussian or CountSketch projection would apply only to ridge; we did not run it, see Section~8.)

%-----------------------------------------------------------------------
\section{Time (Plot 2) and the speed-up table}

$T_\text{build}$ includes fitting the five rough models on 20k-row uniform samples and scoring all \n{1,133,640} training
rows; $T_\text{train}$ is the weighted fit on the summary; speed-up $=T_\text{full}/(T_\text{build}+T_\text{train})$.

\begin{figure}[t]
\centering\includegraphics[width=\linewidth]{\FIG fig_plot2_time.pdf}
\caption{\textsf{PLOT 2.} End-to-end time against size (log--log). For LightGBM the coreset's fixed build cost (about 10 s)
dominates below 25k rows; for ridge the build (18.6 s) is larger than the full fit (4.3 s), so the coreset never saves time.}
\end{figure}

{\small
\begin{tabularx}{\linewidth}{@{}r M M M O M@{}}
\toprule
\thh{LightGBM} & \thh{$T_b$ (s)} & \thh{$T_t$ (s)} & \thh{RMSLE} & \thh{speed-up} & \thh{unif. speed-up} \\
\midrule
10k  & 10.2 & 7.4  & 0.1170 & 3.25$\times$ & 8.09$\times$ \\
25k  & 10.2 & 9.5  & 0.1108 & 3.05$\times$ & 6.28$\times$ \\
50k  & 10.2 & 11.4 & 0.1070 & 2.75$\times$ & 5.18$\times$ \\
100k & 10.2 & 14.0 & 0.1058 & 2.45$\times$ & 4.39$\times$ \\
200k & 10.2 & 21.8 & 0.1045 & 1.85$\times$ & 1.63$\times$ \\
\bottomrule
\end{tabularx}}
\captionof{table}{Speed-up against $T_\text{full}=\n{59.7}$ s (LightGBM, 16 threads).}

\textbf{The trade-off.} Uniform is always faster (no build), the coreset is always more accurate from 5k on. To reach a
given error the coreset needs fewer rows: uniform at 100k (\n{0.1099}, 4.39$\times$) is worse than the coreset at 50k
(\n{0.1070}, 2.75$\times$). At 200k rows the coreset is both faster and better than uniform (1.85$\times$ vs 1.63$\times$),
because the weighted sample trains faster than an equally large uniform one in this run.

\begin{limit}{Timings are development timings}
All times were measured on a shared 128-core server with a load average of 10--30 from other users' jobs. Ratios between
methods are stable across our reruns, but the absolute seconds are not final; the exclusive-server rerun (fixed core set,
16 threads, load logged) is listed in Section~8 and has not been done yet.
\end{limit}

%-----------------------------------------------------------------------
\section{The coreset (Plots 3--5)}

\textbf{Rough models.} Five models $\tilde x_1..\tilde x_5$ fitted on 20k-row uniform samples with different seeds. Their
roughness is $\alpha_t=F(P,w,\tilde x_t)/\text{opt}$ with opt $=F(P,w,x^\ast)$; the largest is \n{2.45} for LightGBM and
\n{1.01} for ridge (a 20k-row ridge is already almost optimal).

\textbf{Sensitivities.} With $f_i$ the squared log error of row $i$ and $F_\text{sum}(\tilde x_t)=\sum_j w_jf_j(\tilde x_t)$,
\begin{eqbox}
\[ s_i=\max_{t=1..5}\Big[\frac{f_i(\tilde x_t)}{F_\text{sum}(\tilde x_t)}+\frac1n\Big], \]
\[ S=\sum_{i=1}^n s_i,\qquad p_i=s_i/S,\qquad w_i=\frac{1}{m\,p_i}. \]
\end{eqbox}
Using the sum (not the mean) in the denominator makes $\sum_i f_i/F=1$ per rough model, so $S$ must lie in $[2,6]$ for five
models and one group; a unit test enforces it. Measured: $S=\n{2.47}$ (LightGBM), \n{2.12} (ridge). The rough models score
every training row; no test row is ever scored by them. For trees there is no $(1\pm\varepsilon)$ guarantee; we call the method
\emph{heuristic sensitivity-motivated importance sampling}.

\textbf{Why it should help here.} A row's sensitivity is large when some rough model fits it badly: event days, hubs after a
data gap, sharp promotion jumps. Uniform sampling sees such rows in proportion to their count (a few per cent), while the
coreset over-samples them and down-weights them by $1/(mp_i)$, so the weighted loss stays an unbiased estimate of the full
loss but with lower variance on exactly the rows that drive the error. The $1/n$ term keeps every row reachable, which bounds
the largest weight. The histogram below shows the mass is spread over three orders of magnitude, with a thin tail of rows
up to about 250 times the uniform share $1/n$.

\begin{figure}[h]
\centering\includegraphics[width=\linewidth]{\FIG fig_hist_sensitivity.png}
\caption{\textsf{SENSITIVITY SCORES} $s_i/S$ over all \n{1,133,640} training rows (log axes). The floor at
$10^{-6.4}\approx1/(nS)$ is the $1/n$ term; the tail is what the coreset over-samples.}
\end{figure}

\begin{figure*}[t]
\centering\includegraphics[width=\linewidth]{\FIG fig_plot345_lgbm.pdf}
\caption{\textsf{PLOTS 3--5 (LightGBM).} Left: relative training-loss gap; the coreset is lower at every size (\n{0.373} vs
\n{0.557} at 50k). Middle: test loss over the full model's \emph{training} loss; it levels off near the full model's own
generalisation ratio $F(P_\text{test},x^\ast)/F(P,x^\ast)-1=\n{1.17}$, not at 0. Right: relative distance between the summary
model's and $x^\ast$'s test predictions; the coreset is closer at every size (\n{0.0049} vs \n{0.0056} at 50k).}
\end{figure*}

\textbf{Slopes (log--log, LightGBM).} Plot 3: coreset \n{-0.59}, uniform \n{-0.48}; Plot 5: \n{-0.19} and \n{-0.18};
ratio Plot~5/Plot~3 for the coreset \n{0.33}. \textbf{Ridge:} Plot 3 \n{-0.78} / \n{-0.99}; Plot 5 \n{-0.52} / \n{-0.50};
ratio \n{0.66}. For ridge the parameter error should fall at about half the rate of the loss gap (a gap $\varepsilon$ in the
loss corresponds to $\sqrt{\varepsilon}$ in the parameters); the measured \n{0.66} is close. For trees no such relation is
guaranteed and the ratio is smaller (\n{0.33}): the prediction vector moves slowly even when the training loss gap closes fast.

\textbf{Ridge.} The coreset helps the training gap at small sizes (\n{0.030} vs \n{0.084} at 2k), but the test error is
the same for every method (\n{0.159}--\n{0.166}) and the build costs more than the full fit. Honest negative: for a cheap
convex model, sampling saves nothing.

\begin{figure*}[t]
\centering\includegraphics[width=\linewidth]{\FIG fig_plot345_ridge.pdf}
\caption{\textsf{PLOTS 3--5 (RIDGE).} Here the theory is exact. The coreset closes the training gap faster at small sizes
(left), the test ratio is flat for both methods (middle: the model is too simple to over-fit), and the coefficient distance
(right) falls at \n{0.66} times the rate of the training gap, close to the predicted one half.}
\end{figure*}

\textbf{Plot 4 needs care.} Its denominator is the \emph{training} loss of $x^\ast$, so even $y^\ast=x^\ast$ gives
$F(P_\text{test},x^\ast)/F(P,x^\ast)-1=\n{1.17}$ for LightGBM: the full tree model fits its training rows about twice as well
as unseen days. Both methods therefore level off near \n{1.17}--\n{1.3}, and the coreset approaches that floor faster
(\n{1.38} vs \n{1.52} at 50k). For ridge the floor is about \n{0.16}, and the curves sit on it from 5k rows on.

\begin{stamp}{Checks that held exactly}
$S$ in $[2,6]$ for both models (\n{2.47}, \n{2.12}) \textbullet\ 5 repeats $\times$ 7 sizes $\times$ 3 methods $\times$ 2 models
= \n{210} summary fits plus \n{50} column fits, all in \code{results/endterm\_run1\_results.csv} \textbullet\ same model settings on
full data and every summary \textbullet\ test rows never used for fitting, sampling probabilities or column choice.
\end{stamp}

%-----------------------------------------------------------------------
\section{Track B, the forecaster we submit: what we tried, what failed and why}

Separately from the coreset study we improved the forecaster whose output is the submission file. It starts from the same
LightGBM recipe as $x^\ast$; the coreset study keeps \code{cs02} as $x^\ast$ (reading guide, Section~1). Evaluation uses
rolling forecast origins: development folds 1--4 start on 2 Jan, 13 Feb, 27 Mar and 8 May 2015 (46{,}830 rows each); fold 5
starts on 19 Jun 2014 and is season-matched to the real forecast window, which also has no regional holidays.

\begin{eqbox}\small\textbf{Adoption gate, fixed before the runs.} A change is adopted only if (a) its mean gain over
the baseline on folds 2--4 is $\ge\n{0.0024}$, twice the measured seed-to-seed noise \n{0.0012}; (b) at least 2 of the 3
folds improve; (c) the sign is the same for seeds 20260807/08/09; (d) it still holds without open-but-zero rows; (e) fold 5
is checked. Better on most folds but missing the gate = \emph{challenger}, not adopted.
\end{eqbox}

{\small
\begin{tabularx}{\linewidth}{@{}Y M@{}}
\toprule
\thh{Step (mean RMSLE, folds 1--4)} & \thh{RMSLE} \\
\midrule
Earlier 3-model blend (CatBoost + LightGBM + Tweedie), re-scored on corrected folds & 0.1250 \\
LightGBM \code{base} (feature fraction 0.8) & 0.1312 \\
\code{cs02}: feature fraction 0.2 + row bagging, 1 seed (seeds 0.1190 / 0.1186 / 0.1191) & 0.1190 \\
\code{cs02}, 3 seeds averaged in log space & 0.1173 \\
\code{cs02}, lr 0.025, 1000 fixed rounds, 3 seeds (\emph{like-for-like baseline}) & 0.1170 \\
\code{hp17}, same setting, 3 seeds & 0.1141 \\
\textbf{\code{hp17} + TabPFN-TS (mix of contexts), weight 0.2 (frozen)} & \textcolor{ultra}{\textbf{0.1129}} \\
\bottomrule
\end{tabularx}}
\captionof{table}{From the earlier blend to the frozen model. Last two rows: our re-computation from saved predictions.}

\textbf{What worked, and why.} \emph{Feature fraction 0.2 with bagging} ($0.1312\to0.1190$): with 342 correlated
columns (many rolling windows of the same series), each tree sees a small random subset, which decorrelates the trees and
stops a few long-window medians from dominating. \emph{Seed averaging} (\n{-0.0017}): free variance reduction.
\emph{\code{hp17}} (255 leaves, min child 160, L2 24.1, feature fraction 0.15): bigger trees with strong shrinkage of leaf
values fit hub-specific weekday and promotion patterns without over-fitting small hubs. \emph{TabPFN-TS}: a pretrained
tabular foundation model run per hub on its own history; its errors are different from the trees' (it reads each series'
recent level directly), so a 20\% log-space blend helps.

\textbf{Searches (all on the same folds, so the best is optimistic).} A 20-config random search found \code{hp17} as the only
config better on all four folds (best of 19). Eight neighbours (L2 10/50, 127/511 leaves, min child 320, feature fraction
0.1/0.2, fixed rounds) all landed within \n{0.002}: a plateau, not a lucky point. A pooled Optuna search (TPE, 2 sampler
seeds, 80 trials, objective = folds 2--3 blend RMSLE) found \n{0.1105} against \n{0.1109} for \code{hp17}: \n{0.0004}, inside the
noise, so no replacement; its five finalists, re-run on all five folds, gained only \n{+0.0001} to \n{+0.0007} over
\code{hp17} on folds 2--4 (gate \n{0.0024}). TabPFN-TS context lengths 365 / 540 / 730 days and weights 0.1--0.5 were compared offline; the grid is
flat within \n{0.001}--\n{0.002}.

\begin{figure}[h]
\centering\includegraphics[width=0.92\linewidth]{\FIG fig_optuna.pdf}
\caption{\textsf{OPTUNA, 80 TRIALS.} The ten best trials span only \n{0.0005}; the whole good region (255--383 leaves,
feature fraction 0.17--0.21) sits inside one noise band. Searching further would add selection bias, not accuracy.}
\end{figure}

{\small
\begin{tabularx}{\linewidth}{@{}l M M M M M M@{}}
\toprule
\thh{config} & \thh{col} & \thh{leaves} & \thh{minc} & \thh{L2} & \thh{4-fold} & \thh{gain 2--4} \\
\midrule
\textbf{hp\_017} & 0.15 & 255 & 160 & 24.1 & \textcolor{ultra}{\textbf{0.1149}} & +.0025 \\
hp\_006 & 0.10 & 63 & 320 & 22.7 & 0.1174 & +.0002 \\
hp\_014 & 0.15 & 255 & 160 & 2.9 & 0.1175 & $-$.0002 \\
hp\_008 & 0.10 & 255 & 40 & 0.7 & 0.1178 & +.0007 \\
hp\_000 (=\code{cs02}) & 0.20 & 63 & 80 & 1.0 & 0.1190 & 0 \\
hp\_016 & 0.15 & 31 & 80 & 3.6 & 0.1234 & \textcolor{amber}{$-$.0027} \\
hp\_004 & 0.15 & 31 & 160 & 0.35 & 0.1242 & \textcolor{amber}{$-$.0066} \\
hp\_005 & 0.40 & 31 & 320 & 7.1 & 0.1252 & \textcolor{amber}{$-$.0060} \\
\bottomrule
\end{tabularx}}
\captionof{table}{20-config random search (seed 07, lr 0.05, early stopping): top 4, control, bottom 3. Pattern: 255 leaves +
strong L2 + large min child win; every 31-leaf tree loses on every fold.}

\textbf{TabPFN-TS: context length and weight.} The member reads the last $C$ days of each hub. With $C=365$ it loses
\n{0.0024} on winter fold~1 (our reading: with one year of context it sees a single year-end and copies it); $C=730$ fixes fold~1 but loses on
fold~5, where only about 535 days of history exist. The equal log-space mix of both is the only input that does not lose on
fold~1 and fold~5 at the same time, and its folds 2--4 gain (\n{0.0045}) is within noise of 730 alone (\n{0.0053}).

\begin{figure*}[t]
\centering\includegraphics[width=\linewidth]{\FIG fig_tabpfn_grid.pdf}
\caption{\textsf{TABPFN-TS GRID} (offline on saved predictions; \code{hp17} 3-seed average + TabPFN-TS at weight $w$, gain over
\code{cs02}). Look at the dotted line ($w=0.2$): the mix (blue) is never the best curve on any single panel, but it is the only
one that is positive and near the top on all three. Higher weights help folds 2--4 by $\le$\n{0.0010} (below noise) and hurt fold~5.}
\end{figure*}

{\footnotesize
\begin{tabularx}{\linewidth}{@{}L{2.55cm} L{2.15cm} Y@{}}
\toprule
\thh{Idea} & \thh{Measured} & \thh{Why it failed (our reading)} \\
\midrule
Level target, offsets, recency weights & lose or tie & trees already learn level from the rolling medians; a global shift cannot fix hub-specific bias \\
Tweedie, CatBoost, 3-model blend & 0.1178 vs 0.1173 & members too correlated with LightGBM; blend adds nothing over seed averaging \\
NNLS blend weights; calibration $a+b\hat y$ & 0.1192 vs 0.1184; worse & weights fitted on few folds over-fit; calibration is unstable across seasons \\
Horizon split, per-format models & 0.1190; worse & smaller training sets per model lose more than they gain \\
New features (trend, streaks, year-ago, peers; 3 rounds) & fold 4 $-$0.002..$-$0.005, fold 3 flips & lead did not replicate on new seeds; signal already in existing columns \\
Event-day features & worse or equal & the target window has no regional holidays \\
MLP, Chronos-2, TiDE, NHITS & best blend +0.0019 on one fold & weaker alone (MLP fold 1 0.2454); gains below noise \\
Summer origins, 7-day origins & +0.0013; worse & noise; doubles memory \\
SVD, two-stage, peer residuals, reconciliation, regime experts & $\le$ +0.0027 on 1 of 3 folds & below gate or fold-specific \\
Shift adaptation & +0.0011..0.0016, all folds & consistent but below gate; needs an extra forecast \\
Importance pruning (81 cols) & $-$0.0063 / $-$0.0051 & low-importance columns still help on some folds \\
Per-tier member on \code{hp17} & worse than plain & tier split shrinks the data for big trees \\
Hurdle model & $\le$1 open-zero row per fold & nothing to model \\
\bottomrule
\end{tabularx}}
\captionof{table}{Experiment ledger, negatives included (more than 40 ideas; full table in \code{results/experiment\_ledger.md}).}

Error analysis of \code{cs02}: the 10 worst of 168 forecast days carry 28\% of the squared error, and they are event days
(holidays, Easter, Sundays, school closures; RMSLE \n{0.267} vs \n{0.133}). This is why feature and model ideas that target
ordinary days gave so little.

%-----------------------------------------------------------------------
\section{Freeze audit: independent re-computation}

Before freezing, every gate number was recomputed from the saved out-of-fold prediction files with a separate script
(\code{code/freeze\_audit.py}), not taken from the run logs. Candidate formula, for hub $h$ and day $t$
($\ell=\log(1+\cdot)$, closed days $=0$, clipped at 0):
\begin{eqbox}\small
\[ L_{h,t}=0.8\cdot\tfrac13\!\!\sum_{s\in\{07,08,09\}}\!\!\ell\big(\hat y^{\,\text{hp17}}_{s}\big) \]
\[ \qquad +0.2\cdot\tfrac12\Big(\ell\big(\hat y^{\,\text{TPF},365}\big)+\ell\big(\hat y^{\,\text{TPF},730}\big)\Big) \]
\[ \hat y_{h,t}=\max\big(0,\;e^{L_{h,t}}-1\big)\ \text{if open, else }0. \]
\end{eqbox}

\begin{figure*}[t]
\centering\includegraphics[width=0.82\linewidth]{\FIG fig_gate_gains.pdf}
\caption{\textsf{GATE.} Gain over the like-for-like baseline. \code{hp17} alone misses the gate on folds 2 and 4; with
TabPFN-TS (mix of contexts, blue) all three folds 2--4 clear it. The 365-day-only fallback loses on winter fold~1 (hatched).}
\end{figure*}

{\small
\begin{tabularx}{\linewidth}{@{}Y M M M M M M@{}}
\toprule
\thh{gain vs cs02} & \thh{f1} & \thh{f2} & \thh{f3} & \thh{f4} & \thh{f5} & \thh{2--4} \\
\midrule
seed 07 & +.0024 & +.0055 & +.0044 & +.0043 & +.0035 & .0047 \\
seed 08 & +.0038 & +.0043 & +.0062 & +.0034 & +.0075 & .0047 \\
seed 09 & +.0039 & +.0040 & +.0065 & +.0032 & +.0067 & .0046 \\
3-seed avg & +.0031 & +.0045 & +.0054 & +.0035 & +.0057 & \textcolor{ultra}{\textbf{.0045}} \\
excl. open-0 & & & & & +.0089 & .0046 \\
\midrule
\code{hp17} alone & +.0058 & +.0005 & +.0045 & +.0007 & +.0038 & \textcolor{amber}{.0019} \\
\bottomrule
\end{tabularx}}
\captionof{table}{Recomputed gate (fold 5 = the 935 hubs with a full 42-day block, 39{,}270 rows).}

Absolute 3-seed RMSLE, baseline $\to$ frozen model: fold 1 \n{0.1334$\to$0.1302}, fold 2 \n{0.1186$\to$0.1141}, fold 3
\n{0.1128$\to$0.1073}, fold 4 \n{0.1034$\to$0.0998}, fold 5 \n{0.1794$\to$0.1737}. All five gate conditions pass on every seed;
\code{hp17} alone fails (a) and (d).

{\footnotesize
\begin{tabularx}{\linewidth}{@{}L{0.35cm} L{2.4cm} L{0.85cm} Y@{}}
\toprule
\thh{\#} & \thh{Check} & \thh{Result} & \thh{Evidence} \\
\midrule
1 & Formula = implementation & \key{PASS} & log1p blend 0.8/0.2, closed 0, clip 0, \code{hp17\_fix1000} registry check; rehearsal output rebuilt from members with max $|\Delta\ell|=0$ \\
2 & No leakage & \key{PASS} & training labels end 19 Jun 2015 = forecast origin; TabPFN history ends 19 Jun (both manifests); hub encoding cross-fitted on training rows only; no test labels exist \\
3 & Gate from OOF files & \key{PASS} & table above; matches the run logs within \n{0.0001} rounding \\
4 & 365/730 mix & \key{PASS} & rebuilt from the two files: max $|\Delta\ell|=0$ on 187{,}320 dev rows and 46{,}830 test rows; SHA-256 matches \\
5 & Bagging active & \key{PASS} & saved boosters: \code{bagging\_fraction 0.85}, \code{bagging\_freq 1}, 1000 trees \\
6 & Final retrain & \warn{COND.} & real path, own folder, nothing overwritten; 13/13 sanity checks pass; a second run of the same recipe is not bit-identical (RMS log difference \n{0.0017}; seed-to-seed is \n{0.019}) \\
7 & Wording & \key{PASS} & no outside scores, rankings or ``state of the art''; folds 1, 4, 5 called diagnostics \\
\bottomrule
\end{tabularx}}
\captionof{table}{Freeze audit (full write-up: \code{docs/FREEZE\_AUDIT.md}).}

\begin{stamp}{Frozen artifact (final retrain, 26 Nov 2026) \textendash\ decision: GO}
Recipe \n{hp17\_fix1000}, 3 seeds, 37 origins (labels end 2015-06-19), \n{1,434,050} rows $\times$ \n{342} columns; TabPFN-TS mix at
$w=0.2$ in log space \textbullet\ boosters SHA-256: \n{613aa09d\ldots a816}, \n{d0a4af42\ldots dc4a}, \n{fa798ad0\ldots 6ba7}
\textbullet\ blend file \n{6d5eb125\ldots cab3} (= rehearsal 2) \textbullet\ forecast file \n{7d44a050\ldots 4484}: \n{46,830} rows,
Id order = sample file, 0 NaN, 0 negative, \n{6,548} closed rows exactly 0, 0 open rows at 0; formula rebuilt from the saved
members with max $|\Delta\ell|=\n{0.0}$; 13 of 13 sanity checks pass. Against rehearsal 2 (same recipe, same inputs): RMS log
difference \n{0.0017}, correlation \n{0.99999} \textendash\ not bit-identical (multi-threaded LightGBM on a loaded server; see the
audit). Fallbacks kept unchanged: rehearsal 2, rehearsal 1 (365-day input), \code{cs02}.

\end{stamp}

\begin{limit}{Selection effects we cannot remove}
\code{hp17} was the best of 19 configs and its gain shrank on fresh seeds; the 365/730 mix was defined after seeing both;
the weight 0.2 was the middle of a flat curve on folds 2--4, which include fold 4; TabPFN-TS was configured on a 100-hub fold-4
smoke run. About 120 recipe variants were scored on the same folds (20 + 8 + 80 + contexts + weights + about 40 ideas). Folds 1,
4 and 5 are confirmation diagnostics, not untouched test sets, and the real forecast window has no labels in this project, so no
real test score exists.
\end{limit}

%-----------------------------------------------------------------------
\section{Reproducibility, and questions we expect}

\textbf{What is in the repository} (\code{3\_endsem/}): \code{notebook/01\_endterm\_results.ipynb} (executed; rebuilds every
table and figure of this report from \code{results/}, recomputes the gate from the saved out-of-fold files and runs the
leakage check, in under a minute); \code{notebook/02\_run\_frozen\_forecaster.ipynb} (loads the three frozen boosters and their
exact input matrix, \code{final/forecast\_matrix\_open.parquet}, and rebuilds the submitted forecast file with a maximum
log difference of \n{0.0} in about 10 seconds; both notebooks clone the repository themselves when opened on a fresh machine
or Colab); \code{code/} (end-term experiment, coreset and plotting scripts, the final-retrain and sanity-check scripts, the
audit, the leakage check); \code{final/} (the three frozen LightGBM boosters, compressed, their input matrix, the forecast file
and the manifest with SHA-256 of every file);
\code{artifacts/} (out-of-fold predictions, TabPFN-TS predictions, both rehearsal runs); \code{MANIFEST.sha256} (SHA-256 of
every file). The heavy steps (building 37 origin tables, 260 end-term fits, TabPFN-TS on a GPU) need the full data and
1--2 hours on a server; the notebook runs on a laptop from the saved outputs.

{\footnotesize
\begin{tabularx}{\linewidth}{@{}L{2.6cm} Y@{}}
\toprule
\thh{Question} & \thh{Answer from our files} \\
\midrule
Why sensitivity sampling for trees? & No guarantee exists, but rows that rough models fit badly are the ones a small sample
  misses; Plot~1 shows the empirical gain (\n{0.0032} at 50k) and Plot~5 shows the summary model stays closer to $x^\ast$. \\
Why ridge at all? & It is the case where the theory is exact, so it checks our implementation: $S=\n{2.12}$ and the
  Plot~5/Plot~3 slope ratio \n{0.66} behave as predicted. \\
What is the noise floor? & Seed-to-seed change of the same recipe: \n{0.0012} on the 3-fold mean. The gate is twice that. \\
Is fold~4 a clean test? & No: it is also development fold~4. Removing it from every decision changes none of them
  (Section~1), but we still report the test number as a development-fold result. \\
Which model is submitted? & \code{hp17} (3 seeds) + TabPFN-TS at 0.2 in log space; its window has no labels, so it has
  development-fold scores only. The coreset study uses \code{cs02} as $x^\ast$. \\
Why not coreset the blend? & TabPFN-TS has no training loss or row weights (pretrained, reads history at prediction time), so
  sensitivities are undefined; the LightGBM part could be, at several hours of server time, and is not needed for the rubric. \\
Which features matter? & The hub $\times$ weekday $\times$ promotion profile (\n{46\%} of split gain) and recent rolling levels
  (\n{22\%}); 115 session columns add \n{2.5\%}. \\
Why not tune more? & 80 Optuna trials found \n{0.0004} over \code{hp17}, a third of the noise; more trials add selection bias. \\
Does the time split shift? & Yes: winter fold~1 behaves differently (TabPFN-TS 365 loses there); the forecast window is summer, matched by fold~5. \\
\bottomrule
\end{tabularx}}

%-----------------------------------------------------------------------
\section{Limitations and open items}

\begin{itemize}
\item \textbf{Timing.} Development timings only; exclusive rerun pending (Section~4).
\item \textbf{Fold 4 = test window.} No decision changes without it (Section~1), but the test RMSLE of $x^\ast$ is not from an
  untouched holdout; no labelled window was left unused.
\item \textbf{Winter.} The frozen model is \n{0.0027} worse than \code{hp17} alone on fold~1 (winter, year-end); the forecast
  window is summer and fold~5 improves, so we accept it, but it is unexplained.
\item \textbf{Trees.} The coreset has no guarantee for trees; its advantage at mid sizes is empirical (5 repeats, 1 split). A
  check with unnormalised weights at one size was planned and not run.
\item \textbf{Columns.} Gaussian/CountSketch projections for ridge were not run; only the column-selection experiment for trees.
\item \textbf{Fold 5} covers 935 of 1{,}115 hubs (180 hubs lack 2014 H2 data); its error is larger because of 9 open-zero rows.
\item \textbf{Closed rows} inflate $R^2$; read RMSLE first.
\end{itemize}

%-----------------------------------------------------------------------
\section{Conclusion}

\textbf{Chosen summary size: $m=50{,}000$ rows with the sensitivity coreset.} There the LightGBM test RMSLE rises from
\n{0.1022} to \n{0.1070} (\n{+0.0048}), RMSE on the log scale is the same number, and $R^2$ stays above \n{0.998}; end-to-end
time falls from \n{59.7} s to \n{21.6} s (\n{2.75$\times$}). Uniform sampling at the same size is faster (\n{5.18$\times$}) but
\n{0.0032} worse, and needs two to four times as many rows to match. For ridge, sampling saves no time and changes nothing:
the build costs more than the full fit. Separately, the forecaster we submit (\code{hp17} + TabPFN-TS, 0.8/0.2 in log space)
improves on the \code{cs02} recipe by \n{0.0045} RMSLE on development folds 2--4, passed a gate fixed in advance (also with
fold~4 removed), and was frozen with hashes, sanity checks and a fallback; its forecast window has no labels, so it has no
test score.

{\footnotesize\textbf{References.} [1] O.~Bachem, M.~Lucic, A.~Krause, \emph{Practical coreset constructions for machine learning},
arXiv:1703.06476, 2017. [2] D.~P.~Woodruff, \emph{Sketching as a tool for numerical linear algebra}, arXiv:1411.4357, 2014.
[3] P.~Ma, M.~W.~Mahoney, B.~Yu, \emph{A statistical perspective on algorithmic leveraging}, arXiv:1306.5362, 2013.
[4] S.~Hoo et al., \emph{From tables to time: how TabPFN-v2 outperforms specialized time series forecasting models},
arXiv:2501.02945, 2025. [5] G.~Ke et al., \emph{LightGBM: a highly efficient gradient boosting decision tree}, NeurIPS 2017.
[6] CSE558 lectures 13--22 (sensitivity sampling, Johnson--Lindenstrauss, sketching for least squares).\par}

\end{document}
